% =============================================================================
% 1. ABSTRACT
% =============================================================================
\vspace{-2.5ex}

\begin{abstract}
Coding agents must integrate external tool returns into ongoing reasoning---a
capability that standard left-to-right pretraining on code exposes only in its
forward direction. We observe that the
\textit{action $\rightarrow$ observation $\rightarrow$ continuation} loop of a
coding agent is structurally isomorphic to a function call site, where a
caller binds arguments, a callee returns a value computed elsewhere, and
downstream code consumes that value. This conditioning structure exists at
internet scale in ordinary code. We exploit it through
\emph{function-aware fill-in-the-middle (FIM) mid-training}: a self-supervised
objective that masks functions selected via program dependency graph
analysis and a complexity--inferability double criterion. We mid-train
Qwen2.5-Coder (7B/14B/32B) and Qwen3-8B on a 3B-token decontaminated corpus
drawn from 990 GitHub repositories, then apply existing agentic post-training
pipelines (R2E-Gym, SWE-Smith) without modification. Mid-training improves
SWE-Bench-Verified by $+2.8/+3.0/+3.3$ points at 7B/14B/32B---a consistent
gain across scales---and the improvement holds across post-training pipelines
and base-model families. Beyond in-domain gains, mid-training also mitigates
the capability erosion that agentic post-training otherwise inflicts on
non-agent coding (e.g., LiveCodeBench) and non-coding tool-use benchmarks
(tau-bench, BFCL): although the mid-training corpus contains Python code
only, the function-call inductive bias survives post-training and yields
consistent improvements over post-training alone.
\end{abstract}

\vspace{-1.5ex}